\documentstyle[%


righttag%


,11pt%


,amssymb%


,russian%               remove in Eng.


,izvru%                 Set izven% in Eng.


]{amsart}





% {,} !!


% repl. \rule with smth. else


\newcommand{\rank}{\operatorname{rank}}


\newcommand{\tts}[1]{}





%\hoffset=-15mm%                  For Eng.


\setcounter{page}{59}


\god{2004}


\nomer{6~(505)} %                        Remove in Eng.


%\nomer{Vol.\,48, No.\,6}%               Set in Eng.


%\str{\pageref{begin}--\pageref{end}}%  Set in Eng.


\udk{519.929}





\author{\it В.Б.\,Черепенников}


% NB:   ^^ remove in Eng.


\title{Обратная начальная задача для


некоторых линейных систем


дифференциально-разностных уравнений}


\thanks{Работа выполнена при частичной финансовой поддержке Российского


фонда фундаментальных исследований, грант \No\,03-01-00203.}





\date{}


%\pagestyle{myheadings}%         Set in Eng.


\pagestyle{plain} %              Remove in Eng.


\begin{document}


%\thispagestyle{plain}%          Set in Eng.


\maketitle


\label{begin}








В числе наиболее исследованных систем


функционально-дифференциальных уравнений находятся линейные системы


дифференциально-разностных уравнений (ЛС ДРУ), когда запаздывание


постоянно ([1]--[5]).


Во многих работах в качестве основной рассматривалась


начальная задача, когда на начальном множестве тем или иным способом


задавалась начальная функция. Для ЛС ДРУ было показано, что задание


начальной функции гарантировало существование единственного решения


как в положительном направлении оси независимой переменной, так и в


отрицательном. В ряде случаев возникала


необходимость найти решение, удовлетворяющее некоторым дополнительным


(напр., краевым) условиям. Изучению вопросов разрешимости


краевых задач для различных типов функционально-дифференциальных


уравнений посвящено значительное количество работ (см. библиографию в


[5]).





В последнее время в некоторых прикладных областях науки сформировалась


задача для ЛС ДРУ, связанная с исследованиями начальной функции


[6]. Назовем ее обратной начальной задачей для ЛС ДРУ. Она


формулируется следующим образом: найти условия существования 


начальной абсолютно непрерывной функции такой, что порождаемое ею 


решение начальной задачи удовлетворяет дополнительным известным 


условиям, наложенным на решение и (или) его производные.


В ряде работ (напр., [7], [8], [5])


аналогичная задача для определенных краевых условий формулируется как


задача управления.





Исследованию обратной начальной задачи для линейных систем


дифференциально-раз\-ност\-ных уравнений посвящена данная статья.








\section{Постановка задачи}





Рассмотрим начальную задачу для линейной системы


дифференциально-разностных уравнений запаздывающего типа


\begin{gather}


\Dot{x}(t)=A(t)x(t-1)+B(t)x(t)+f(t),\ \ t\in J=[0,T];\\


%


x(t)=g(t),\ \ t\in J_0=[-1,0],


\end{gather}


где $x(t):J_0\cup J \to R^n$; $A(t), B(t):J \to R^{n\times n}$;


$g(t):J_0 \to R^n$; $f(t):J \to R^n$.





Следуя [3], введем 





\begin{defn}


Вещественная матрица


$A(t)=\{a_{ij}(t)\}$, $i=\overline{1,n}$, $j=\overline{1,m}$,


принадлежит классу $C^k$ на $J^*=[t_1,t_2]$ (будем писать


$A(t)\in C^k[J^*]$), если функции $\{a_{ij}(t)\}$,


$i=\overline{1,n}$, $j=\overline{1,m}$, принадлежат классу


$C^k$ на $(t_1,t_2)$ и имеют в точках $t_1$ и $t_2$


соответственно левые и правые производные порядка $k$ и если


функции $a_{ij}^{(k)}(t)$, определенные таким образом на $t_1\leq t


\leq t_2$, непрерывны справа в точке $t_1$ и слева в точке $t_2$.


\end{defn}





С учетом данного определения под решением задачи (1), (2)


понимается функция $x(t)\in C^1[J]$, удовлетворяющая


линейной системе дифференциально-разностных уравнений


в (1) и совпадающая с начальной функцией


$g(t)\in C^0[J_0]$ на $J_0$ [4].





Установление условий существования и способов нахождения единственного 


абсолютно непрерывного решения $x(t)$ по заданной начальной функции 


$g(t)$ составляет прямую начальную задачу для линейной системы


дифференциально-разностных уравнений (1), (2).





Дадим теперь определение обратной начальной задачи для линейной


системы диф\-фе\-рен\-ци\-аль\-но-разностных уравнений (1). Пусть


известны значения функции $x(t)$ и (или) ее производных


$x^{(n)}(t)$, $n=\overline{1,p}$, в некоторых фиксированных точках


$t\in J_0 \cup J$.


\begin{equation}


\begin{split}


&x(t_i^0)=x_i^0,\ \ i=\overline{1,m};\\


&\Dot{x}(t_j^1)=x_j^1,\ \ j=\overline{1,l};\\


&\dotsc\dotsc\dotsc\dotsc\dotsc\dotsc\dotsc\dotsc\\


&x^{(p)}(t_k^p)=x_k^p,\ \ k=\overline{1,r}.


\end{split}


\end{equation}





\begin{note}


В дальнейшем будем считать, что в конечных


последовательностях $\{t_i^0\}_{i=1}^{m}$,


%%%%%%% join in eng.                    ^^^


$\{t_j^1\}_{j=1}^{l},\dotsc, \{t_k^p\}_{k=1}^{r}$ выполняются


условия упорядоченности, т.\,е.


$t_j^i<t_{j+1}^i$, $i=\overline{0,p}$. При этом случаи, когда


$t_{\gamma}^{\alpha}=t_{\delta}^{\beta}$ ($\alpha \neq \beta$), не


исключаются.


\end{note}





Отметим, что в зависимости от конкретных значений точек $t_s^k$ 


данные условия могут включать, в частности, начальные, 


начально-краевые или краевые условия.





\begin{defn}


Задачу нахождения условий существования


и способов вычисления такой начальной функции $g(t)$, $t\in J_0$,


что решение $x(t)$ начальной задачи


(1), (2) удовлетворяет условиям (3), будем


называть {\em обратной начальной задачей} для данной линейной системы


дифференциально-разностных уравнений.


\end{defn}





В связи с определением 2 следует выяснить: при каких условиях


существует такая начальная функция и какими


свойствами она обладает?








\section{Основные результаты}





Перепишем задачу (1), (2) с учетом (3) в виде


обратной начальной задачи


\begin{gather}


\Dot{x}(t)=A(t)x(t-1)+B(t)x(t)+f(t),\ \ t\in J=[0,T]; \\


%


x(t)=g(t),\ \ t\in J_0=[-1,0],\\


%


x(t_i^0)=x_i^0,\ \ i=\overline{1,m};\ \


\Dot{x}(t_j^1)=x_j^1,\ \ j=\overline{1,l}; \dotsc;\ \


x^{(p)}(t_k^p)=x_i^p,\ \ k=\overline{1,r};\\


%


t_j^i \in J_0 \cup J:i=0,\ \ j=\overline{1,m};\ \


i=1,\ \ j=\overline{1,m};\dotsc;\ \ i=p,\;j=\overline{1,r};\notag \\


%


T=\max\{1,t_m^0,t_l^1,\dotsc,t_r^p\}.


\end{gather}





Представим промежуток $J$ в (4) следующим образом:


\begin{equation}


J=\bigcup _{n=1}^{]T[+1}J_n,\quad J_n=


\begin{cases}


[n-1,\ n], & n=\overline{1,]T[};\\


\lbrack \ ] T[,\ T], & n=]T[+1.


\end{cases}


\end{equation}


Здесь $]\cdot[$ --- целая часть числа.





\begin{defn}


Решение $x(t)$ задачи


(4)--(6) на отрезке $J_n$ будем обозначать $x_n(t)$.


\end{defn}





Далее, из всей совокупности фиксированных точек, заданных на оси


$t$ согласно (6), выделим те, которые принадлежат некоторому


отрезку $J_m$, и пусть $\alpha_m$ --- наибольший из верхних


индексов, который встречается среди обозначений, принятых для


выделенных точек.





В дальнейшем будем рассматривать случай, когда на решение $x(t)$


накладываются дополнительные условия


\begin{equation}


x_m(t)\in C^{s_m}[J_m],\ \ s_m\geq \alpha_m.


\end{equation}





Пусть в (6) $\chi$ --- наибольший порядок производной функции


$x(t)$ в точках $t\in J$, т.\,е.


\begin{equation}


\chi=\max \{i\mid t_j^i \in J\},


\end{equation}


и


\begin{equation}


\chi_0=\max \{i\mid t_j^i \in J_0\}.


\end{equation}





\begin{lem}


Пусть для задачи $(4)$--$(6)$


\begin{itemize}


\item[1.] матрицы $A(t)$, $B(t)$ и вектор $f(t)$ принадлежат


с учетом $(10)$ классу $C^0[J]$, если\linebreak $\chi=0,1$, и


$C^{\chi-1}[J]$, если $\chi\geq 2;$


\item[2.] $ \ov{\sigma} =\max\limits_s\{0,(s-1-]t_1^s[)\}$,


$s=\overline{0,\chi}$.


\end{itemize}





Тогда каждая начальная функция $g(t)$, принадлежащая классу


$C^{\sigma}[J_0]$, где 


$\sigma \geq \max \{ \chi_0,\ov{\sigma}\}$, порождает решение


$x(t)$ такое, что на всех промежутках 


$J_n$, $n=\overline{1,]T[+1}$, оно удовлетворяет требуемым в 


силу $(6)$ условиям дифференцируемости.


\end{lem}





{\bf Доказательство}. 


Рассмотрим начальную задачу (4), (5) и пусть некоторая 


начальная функция $g(t)$ принадлежит классу $C^k[J_0]$. С


учетом условия 1 леммы из известной теоремы о сглаживаемости


решений начальной задачи при возрастании аргумента 


([1], с.\,53) вытекает


$$


x(t)\in C^{k+1}[J_1],x(t)\in C^{k+2}[J_2],\dotsc,


x(t)\in C^{k+]T[+1}[J_{]T[+1}].


$$


Пусть для одного из условий в (6) требуется существование


производной $x^{(s)}(t_1^s)$, где точка 


$t_1^s\in J_{]t_1^s[+1}$. Поскольку $x^{(s)}(t_1^s)$ 


принадлежит  отрезку $]t_1^s[+1$, то начальная функция 


$g(t)\in C^k[J_0]$ на этом отрезке будет


порождать решение $\wt{x}(t)$, принадлежащее не ниже, чем


классу $C^{k+]t_1^s[+1}[J_{]t_1^s[+1}]$. Тогда если 


$k\geq \max\{0,s-]t_1^s[-1\}$, то начальная функция


$g(t)\in C^k[J_0]$ порождает решение 


$\wt{x}(t)\in C^s[J_{]t_1^s[+1}]$. А так как $t_i^s<t_{i+1}^s$,


это означает, что данное условие будет выполняться при всех


$t_i^s\in J$, $i=1,2,\dotsc$ Проводя аналогичные выкладки для


$s=1,2,\dotsc,\chi$ и принимая согласно (11) и условию 2 леммы


$k \geq \max \{ \chi_0,\ov{\sigma}\}$, приходим к утверждению 


леммы.





\begin{cor}


Поскольку в силу определения 1 


производные функции $x(t)$ при целочисленных значениях аргумента $t$ 


односторонние, если в начальной точке $t=0$ имеют место соотношения


$$


g^{(s)}(t)|_{t\to 0_{-}}=


x^{(s)}(t)|_{t\to 0_{+}},\ \ s=\overline{0,m},


$$


то при целочисленных значениях аргумента в силу упомянутой выше


теоремы о сглаживании решений и условия 1 леммы вытекает, что


решение $x(t)$ при $t_*=]t[$, $t\in J$, имеет $k=\min \{m+]t[,\chi\}$


непрерывных производных.


\end{cor}





На некотором промежутке $J^\prime =[a,b]$ рассмотрим


совокупность вещественных матриц $H_i(t): J^\prime \to


R^{n\times n}$, $i=\overline{1,k}$.





\begin{defn}


Вещественные матрицы $H_i(t)$,


$i=\overline{1,k}$, будем называть {\em линейно независимыми}, если


соотношение


$$


\sum_{i=1}^{k}H_i(t)c_i=\Theta\ \; \forall t\in J^\prime ,


$$


где $\Theta \in R^{n\times n}$ --- нулевая матрица, а $c_i


\in R^1$ --- некоторые постоянные, справедливо тогда и только


тогда, когда $c_i=0$, $i=\overline{1,k}$.


\end{defn}





Зададим функцию 


\begin{equation}


x_0(t)=\sum_{u=1}^{G}\Psi_0^u(t)g_u,\ \ t\in J_0.


\end{equation}


Здесь $\Psi_0^u(t): J_0 \to R^{n\times n}$ --- некоторые известные


вещественные матрицы, $g_u$ --- неизвестные постоянные векторы.





\begin{defn}


Линейно независимые матрицы


$\Psi_0^u(t)$, $u=\overline{1,G}$, которые для исследуемой задачи


(4)--(6) с учетом леммы 1 выбираются из класса


$C^{\sigma}[J_0]$, будем называть {\em опорными матрицами} решения


$x(t)$.


\end{defn}





Рассматривая функцию $x_0(t)$ в виде (12) как начальную


функцию для задачи (4)--(6), найдем решение этой задачи на


$J_1=[0,1]$.





Перепишем (4)--(6) следующим образом:


$$


\Dot{x}(t)=B(t)x(t)+A(t)g(t-1)+f(t),\ \ x(0)=x_0,\ \ t\in J_1.


$$


Поскольку эти соотношения представляют собой задачу Коши для


линейной системы обыкновенных дифференциальных уравнений,


решение данной задачи на $J_1$ представляется формулой [3]


\begin{equation}


x(t)=X(t)x(0)+\int_{0}^{t}C(t,s)[A(s)g(s-1)+f(s)]ds,


\end{equation}


где $X(t)$ -- фундаментальная матрица системы


$\Dot{x}(t)=B(t)x(t)$,


а $C(t,s)=X(t)X^{-1}(s)$ --- матрица Коши.





Учитывая (12), преобразуем слагаемые в формуле (13)


\begin{gather*}


X(t)x(0)=X(t)\sum_{u=1}^{G}\Psi_{0}^{u}(0)g_u;\\


%


\int_{0}^{t}C(t,s)[A(s)g(s-1)+f(s)]ds= 


\int_{0}^{t}C(t,s)A(s)\sum_{u=1}^{G}\Psi_{0}^u(s-1)g_uds+


\int_{0}^{t}C(t,s)f(s)ds= \\


%


=\sum_{u=1}^{G}\bigg[\int_{0}^{t}C(t,s)A(s)\Psi_{0}^{u}(s-1)ds\bigg]g_u+


\int_{0}^{t}C(t,s)f(s)ds.


\end{gather*}


Тогда в силу определения 3 и этих соотношений представим


(13) в виде


$$


x_1(t)=x(t)=\sum_{u=1}^{G}\Psi_{1}^{u}(t)g_u+ f_1(t),\ \ t \in J_1,


$$


где


\begin{gather*}


\Psi_{1}^{u}(t)=X(t)\Psi_{0}^{u}(0)+


\int_{0}^{t}C(t,s)A(s)\Psi_{0}^{u}(s-1)ds; \\


f_1(t)=\int_{0}^{t}C(t,s)f(s)ds.


\end{gather*}





Продолжая теперь решение аналогичным образом на $J_2$, $J_3$ 


и т.\,д., для $J_M=[M-1,M]$ имеем


\begin{equation}


x_M(t)=\sum_{u=1}^{G}\Psi_M^u(t)g_u+f_M(t),\;\;t\in J_M.


\end{equation}


Здесь


\begin{gather*}


\Psi_{M}^{u}(t)=X(t)\Psi_{M-1}^{u}(M-1)+


\int_{M-1}^{t}C(t,s)A(s)\Psi_{M-1}^{u}(s-1)ds;\\


f_M(t)=\int_{M-1}^{t}C(t,s)[f_{M-1}(s)+f(s)]ds.


\end{gather*}





\begin{note}


В соответствии с определением 5 решение


$x_{M}(t)$ дифференцируемо на $J_{M}$ согласованное с


условиями (6) число раз, т.е.


$$


x_{M}^{(s)}(t)=\sum_{u=1}^{G}\frac{d^s}{dt^s}\Psi_{M}^u(t)g_u


+f_{M}^{(s)}(t).


$$


\end{note}





Переходим к нахождению начальной функции $g(t)$. С учетом определения


1 и (12) один из подходов к решению обратной задачи


(4)--(6) может быть сформулирован следующим образом:


при каких условиях существует представление $g(t)$ в виде (12)


и как определить неизвестные коэффициенты $g_u$?





Для этой цели воспользуемся условиями (6). Введем функцию


$$


\delta(t)=0\ \  \text{при}\ \ t\leq 0;\quad


\delta(t)=


\begin{cases}


t, & \text{если}\ t=]t[;\\


]t[+1, & \text{если}\ t\neq ]t[,


\end{cases}


\ \ 


\text{при}\ \ t>0.


$$





Поскольку каждая точка $t_s^k\in J_{\delta(t_s^k)}$, согласно


(14) и замечанию 2 из (6) получаем


\begin{align}


x_{\delta(t_i^0)}(t_i^0)


&=\sum_{n=1}^G \Psi_{\delta(t_i^0)}^n(t_i^0)g_n


+f_{\delta(t_i^0)}(t_i^0)=x_i^0,\ \ i=\overline{1,m};\notag \\


%


\Dot{x}_{\delta(t_j^1)}(t_j^1)&=\sum_{n=1}^G \frac{d}{dt}


\Psi_{\delta(t_j^1)}^n(t_j^1)g_n+\Dot{f}_{\delta(t_j^1)}(t_j^1)=


x_j^1,\ \ j=\overline{1,l};\notag \\


%


\dotsc\dotsc\dotsc&\dotsc\dotsc\dotsc\dotsc\dotsc\dotsc\dotsc\dotsc\dotsc\dotsc


\dotsc\dotsc\dotsc\dotsc\notag \\


%


x_{\delta(t_k^p)}^{(p)}(t_k^p)&=\sum_{n=1}^G \frac{d^p}{dt^p}


\Psi_{\delta(t_k^p)}^n(t_k^p)g_n+f_{\delta(t_k^p)}^{(p)}(t_k^p)=


x_k^p,\ \ k=\overline{1,r}. 


\end{align}





\begin{note}


Здесь и далее при $t<0$ функцию $f_0(t)$ полагаем равной нулю.


\end{note}





Обозначая


\begin{align}


A_{i\,s}&=\Psi_{\delta(t_i^0)}^s(t_i^0),\ \


i=\overline{1,m},\ \ s=\overline{1,G};\notag \\


%


A_{(m+j)s}&= \frac{d}{dt}\Psi_{\delta(t_j^1)}^s(t_j^1),\ \


j=\overline{1,l},\ \ s=\overline{1,G};\notag \\


\dotsc\dotsc\dotsc&\dotsc\dotsc\dotsc\dotsc\dotsc\dotsc\dotsc\dotsc\dotsc\dotsc\dotsc\notag \\


%


A_{(m+l+\dotsb+k)s}&=


\frac{d^p}{dt^p}\Psi_{\delta(t_k^p)}^s(t_k^p),\ \


k=\overline{1,r},\ \ s=\overline{1,G};\\ 


%


B_i&=x_i^0- f_{\delta(t_i^0)}(t_i^0),\ \ i=\overline{1,m};\notag \\


%


B_{m+j}&=x_i^1-\Dot{f}_{\delta(t_j^1)}(t_j^1),\ \ 


j=\overline{1,l};\notag \\


\dotsc\dotsc&\dotsc\dotsc\dotsc\dotsc\dotsc\dotsc\dotsc\dotsc\notag \\


%


B_{m+j+\dotsb+k}&=x_k^p-f_{\delta(t_k^p)}^{(p)}(t_k^p),\ \


k=\overline{1,r}, 


\end{align}


перепишем соотношение (15) в виде блочной линейной системы


алгебраических уравнений относительно неизвестных векторов $g_n$:


\begin{equation}


Ag=B,


\end{equation}


где $A=\{ A_{ij}\}$, $i=\overline{1,m+l+\dotsb+r}$,


$j=\overline{1,G}$, --- блочная матрица;


$g=\{g_1, g_2, \dotsc, g_G\}^T$ и $B=\{B_1, B_2, \dotsc,


B_{m+l+\dotsb+r}\}^T$ --- векторы, индекс $T$ означает транспонирование.


Таким образом, решение обратной начальной задачи (4)--(6)


свелось к неоднородной линейной системе алгебраических уравнений.





Обозначим


\begin{equation}


Q=m+l+\dotsb+r.


\end{equation}


Тогда, поскольку матрицы $A_{ij}\in R^{n \times n}$, а $g_n \in


R^n$, матрица $A$ будет прямоугольной размера $nQ\times nG$


(квадратной, если $Q=G$), а векторы $g$ и $B$ будут иметь


соответственно размеры $nG$ и $nQ$.





Анализ линейной системы (18) начнем с выяснения вопроса о


совместности системы. Пусть


$$


\ov{A}=


\begin{vmatrix}


A_{11}  & A_{12} & \dotsb & A_{1G} & B_1 \\


A_{21}  & A_{22} & \dotsb & A_{2G} & B_2 \\


\vdots  & \vdots & \ddots & \vdots & \vdots \\


A_{Q1}  & A_{Q2} & \dotsb & A_{QG} & B_Q 


\end{vmatrix}


$$


будет расширенной матрицей линейной


системы (18). Тогда на основании теоремы Кронекера--Капели 


система линейных уравнений (18) тогда и только тогда совместна, 


когда ранг расширенной матрицы $\ov{A}$ равен рангу матрицы $A$. 





\begin{cor}


Если $\rank \ov{A}>\rank A$, то


обратная начальная задача для линейной системы


дифференциально-разностных уравнений (4)--(6) с выбранными в


(12) опорными матрицами $\Psi^n_0(t)$, $n=\overline{1,G}$, неразрешима.


\end{cor}





Такой случай может иметь место, когда число линейно независимых


опорных матриц $\Psi^n_0(t),\;t\in J_0$, меньше числа заданных


согласно (6) условий. С другой стороны, несовместность


линейной системы может быть вызвана некорректными условиями


(6), например, если в одной и той же точке заданы два или


более значений решения $x(t)$ или его производных.





Рассмотрим теперь случай, когда линейная система (18)


совместна. 





\begin{thm}


Пусть для задачи $(4)$--$(6)$


справедливы следующие условия\/${:}$


\begin{itemize}


\item[1.] матрицы $A(t)$, $B(t)$ и вектор $f(t)$ удовлетворяют


условию $1$ леммы $1;$


\item[2.] опорные матрицы $\Psi^n_0(t)$, $n=\overline{1,G}$, в


$(9)$ с учетом леммы $1$ принадлежат классу


$C^{\sigma}[J_0];$


\item[3.] линейная система $(18)$, порожденная в силу условий


$(6)$ решением $x(t)$, совместна и, следовательно,


число опорных матриц удовлетворяет неравенству


$G\geq \frac{1}{n}\rank A$.


\end{itemize}





Тогда обратная начальная задача $(4)$--$(6)$ разрешима и


имеет решение в виде начальной функции


\begin{equation}


g(t)=\sum_{u=1}^G\Psi^u_0(t)g_u,\ \ g(t)\in C^{\sigma}[J_0].


\end{equation}


При этом число решений бесконечно, если $G>\frac{1}{n}\rank A$,


и представление решения в виде $(20)$ единственно, если 


$G=\frac{1}{n}\rank A$.


\end{thm}





\begin{pf}


Принимая во внимание (19), рассмотрим


конечное множество линейно независимых опорных матриц


$\{\Psi^n_0(t)\}_{n=1}^G$, где $G\geq Q$. На основании условий


1--3 теоремы, леммы~1 и замечания 1 в силу условий (6)


получаем совместную линейную систему алгебраических уравнений


(18) размера $nQ\times nG$ относительно неизвестных


векторов $g_n$.





Раскрывая блоки $A_{ij}$ матрицы $A$, а также векторы $g$ и $B$,


перепишем (18) в виде


$$


A=


\begin{vmatrix}


a_{11}  & a_{12} & \dotsb & a_{1\ov G}  \\


a_{21}  & a_{22} & \dotsb & a_{2\ov G}  \\


\vdots  & \vdots & \ddots & \vdots   \\


a_{\ov{Q}1}& a_{\ov{Q}2} & \dotsb & a_{\ov Q\ov G} 


\end{vmatrix}


\begin{vmatrix}


\ov{g}_1 \\ \ov{g}_2 \\ \vdots \\ \ov{g}_{\ov{G}} 


\end{vmatrix}


=


\begin{vmatrix}


b_1 \\ b_2 \\ \vdots \\ b_{\ov{G}} 


\end{vmatrix}.


$$


Здесь $\ov{G}=nG$, $\ov{Q}=nQ$.





Обозначим ранг матрицы $A$ через $R$. Известно, что если


$R<\ov Q$, то в системе $ \ov Q-R$ строк


являются линейно зависимыми. Тогда выбираем $R$ линейно


независимых строк, оставляем в системе уравнения, коэффициенты


которых вошли в выбранные строки, и в левой части этих


уравнений оставляем такие $R$ неизвестных коэффициентов


$\ov{g}_n$ (пусть это $\ov{g}_1,\ov{g}_2,\dotsc,


\ov{g}_R$), что определитель коэффициентов при них отличен от


нуля, а остальные неизвестные коэффициенты объявляем свободными


и переносим в правые части уравнений,


\begin{equation}


\begin{cases}


a_{11}\ov{g}_1+a_{12}\ov{g}_2+\dotsb+ a_{1R}\ov{g}_R= b_1-


a_{1(R+1)}\ov{g}_{R+1}-\dotsb - a_{1\ov{G}}\ov{g}_{\ov{G}},\\


a_{21}\ov{g}_1+a_{22}\ov{g}_2+\dotsb+ a_{2R}\ov{g}_R= b_2-


a_{2(R+1)}\ov{g}_{R+1}-\dotsb - a_{2\ov{G}}\ov{g}_{\ov{G}},\\


\dotsc\dotsc\dotsc\dotsc\dotsc\dotsc\dotsc\dotsc\dotsc\dotsc


\dotsc\dotsc\dotsc\dotsc\dotsc\dotsc\dotsc\dotsc\dotsc\dotsc\\


a_{R1}\ov{g}_1+a_{R2}\ov{g}_2+\dotsb+ a_{RR}\ov{g}_R= b_R-


a_{R(R+1)}\ov{g}_{R+1}-\dotsb - a_{R\ov{G}}\ov{g}_{\ov{G}}.


\end{cases}


\end{equation}





Так как в этой системе ранг матрицы равен $R$, то, придавая


для каждого $i$ неизвестным коэффициентам


$\ov{g}_{R+j}^i$, $j=\overline{1,\ov{G}-R}$,


произвольные числовые значения и разрешая систему (21)


относительно коэффициентов


$\ov{g}_k^i$, $k=\overline{1,R}$, 


получаем согласно (20) соответствующее решение


$g^i(t)$. Отсюда вытекает, что множество решений


обратной начальной задачи (4)--(6) бесконечно. С другой


стороны, если в (18) число линейно независимых матриц


$\Psi_0^n(t)$ таково, что $ G=\frac{1}{n}\rank A$, то линейная


система (18) имеет единственное решение. Тогда представление


решения $g(t)$ обратной начальной задачи (4)--(6)


в виде (20) будет единственным.


\end{pf}








\section{Приложения}





Рассмотрим обратную начальную задачу для скалярного


дифференциально-разностного уравнения запаздывающего типа с линейным


коэффициентом


\begin{equation}


\Dot{x}(t)=(a_0+a_1t)x(t-1), \ \ t\in J=[0,T];\ \


x(t)=g(t), \ \ t\in J_0=[-1,0],


\end{equation}


и условиями


\begin{equation}


x(0)=x_0,\ \ x(t_1)=x_1,\ \ x(t_2)=x_2,\ \ x(t_3)=x_3.


\end{equation}


Согласно условиям теоремы 1 выбираем линейно независимые опорные


функции\footnote{В скалярном случае опорные матрицы представляют


собой скалярные функции. Поэтому здесь вместо термина ``опорные


матрицы'' будем использовать термин ``опорные функции''.}


$\psi_0^i(t)=t^i$, $i=\overline{0,3}$, и ищем решение обратной


задачи (22), (23) в виде


\begin{equation}


g(t)=\sum_{i=0}^{3}\psi_0^i(t)g_i=


\sum_{i=0}^{3}g_it^i,\ \ t\in J_0.


\end{equation}


Найдем теперь решение $x(t)$, порождаемое этой начальной


функцией. В силу определения 3, формул (8) и (14) для $t\in J_m$ имеем


\begin{equation}


x_m(t)=\sum_{i=0}^{3}\psi_m^i(t)g_i,


\end{equation}


где


\begin{equation}


\psi_m^i(t)=\psi_{m-1}^i(m-1)+\int_{m-1}^t (a_0+a_1t)


\psi_{m-1}^i(t-1)dt.


\end{equation}


Из последней формулы вытекает, что $\psi_m^i(t)$ будет


полиномом степени $2m+i$. Представим его в виде


\begin{equation}


\psi_m^i(t)=\sum_{n=0}^{2m+i}\psi_{mn}^it^n.


\end{equation}


Тогда


\begin{equation}


\psi_m^i(t-1)=\sum_{n=0}^{2m+i}\psi_{mn}^i(t-1)^n=


\sum_{n=0}^{2m+i}b_{mn}^it^n.


\end{equation}


Здесь


$$


b_{mn}^i=\sum_{j=0}^{2m+i-n}(-1)^jC_{n+j}^n\psi_{mn+j}^i,


$$


а $C_{n+j}^n$ --- биномиальные коэффициенты.





Учитывая (28), преобразуем интеграл в (26)


\begin{equation}


\int_{m-1}^t (a_0+a_1t)\psi_{m-1}^i(t-1)dt =


\int_{m-1}^t (a_0+a_1t)


\sum_{n=0}^{2(m-1)+i}b_{(m-1)n}^it^ndt =


\int_{m-1}^t\bigg(\sum\limits_{n=0}^{2m+i-1}D_{mn}^it^n\bigg)dt.


\end{equation}


В этом выражении при $i=0$


\begin{equation}


D_{10}^0=a_0,\ \ D_{11}^0=a_1,


\end{equation}


а при $i\geq 1$


\begin{equation}


D_{mn}^i=


\begin{cases}


a_0b_{(m-1)0}^i, & n=0;\\


a_0b_{(m-1)n}^i+a_1b_{(m-1)(n-1)}^i, & n=\overline{1,2(m-1)+i};\\


a_1b_{(m-1)(2(m-1)+i)}^i, & n=2m+i-1.


\end{cases}


\end{equation}


Значение интеграла (29) определяется формулой


$$


\int_{m-1}^t\bigg(\sum_{n=0}^{2m+i-1}D_{mn}^it^n\bigg)dt=


-\sum_{n=1}^{2m+i}\frac{(m-1)^n}{n}D_{m(n-1)}^i+


\sum_{n=1}^{2m+i}\frac{1}{n}D_{m(n-1)}^it^n,\ \ m \geq 1.


$$


Подставляя полученное выражение в (26), имеем


$$


\psi_m^i(t)=\psi_{m-1}^i(m-1)


-\sum_{n=1}^{2m+i}\frac{(m-1)^n}{n}D_{m(n-1)}^i+


\sum_{n=1}^{2m+i}\frac{1}{n}D_{m(n-1)}^it^n.


$$


Сравнивая это соотношение с (27), получаем, что в последней


формуле коэффициенты $\psi_{mn}^i$ представляются в виде


\begin{equation}


\psi_{mn}^i=


\begin{cases}


\psi_{m-1}^i(m-1)-\sum\limits_{n=1}^{2m+i}


\frac{(m-1)^n}{n}D_{m(n-1)}^i, & n=0;\\


\frac{1}{n}D_{m(n-1)}^i, & n=\overline{1,2m+i}.


\end{cases}


\end{equation}


Итак, с учетом (7), (8) и (25) данная формула позволяет


последовательно вычислять порождаемое решение $x(t)$ на


$J=\bigcup\limits_{n=1}^{]T[+1}J_n$. Действительно, для $J_1=[0,1]$


согласно (32) имеем


$$ 


\psi_{10}^i=\psi_0^i(0);\ \ \psi_{1n}^i=\frac{1}{n}


D_{1(n-1)}^i,\ \ n=\overline{1,2+i};\ \ i=\overline{0,3}, 


$$


где значение $D_{1(n-1)}^i$, $i=\overline{0,3}$, определяется по


(30) и (31). Тогда в соответствии с (32) вычисляются функции


$\psi_1^i(t)$, $i=\overline{0,3}$, и на основании (25)


находится порождаемое решение


$$ 


x_1(t)=\sum_{i=0}^{3}\psi_1^i(t)g_i,\ \ t\in J_1. 


$$


Отсюда для $t=1$ получаем


$x_1(1)=\sum\limits_{i=0}^{3}\psi_1^i(1)g_i$. 


Переходим на следующий отрезок $J_2=[1,2]$. Здесь


в соответствии с (32)


$$


\psi_{2n}^i=


\begin{cases}


\psi_{1}^i(1)-\sum\limits_{n=1}^{4+i}\frac{1}{n}D_{2(n-1)}^i,


& n=0;\\


\frac{1}{n}D_{2(n-1)}^i, & n=\overline{1,4+i}.


\end{cases}


$$


Далее подобно выше описанному вычисляются функции


$\psi_2^i(t)$, $i=\overline{0,3}$, и порождаемое решение $x_2(t)$


для $t\in J_2$. Таким же образом получаем решения $x_n(t)$


на всех отрезках $J_n$, $n=\overline{3,]T[+1}$.





Переходим к определению неизвестных коэффициентов $g_n$. Отметим,


что в силу (24) и первого из условий в (23) $g_0=x(0)=x_0$. Для


точек $t_1$, $t_2$ и $t_3$ согласно (23) и (15) имеем


$$


x_{\delta(t_s)}(t_s)=\sum_{i=0}^3


\psi_{\delta(t_s)}^i(t_s)g_i=x_s,\ \ s=\overline{1,3}.


$$


Поскольку коэффициент $g_0$ известен, перепишем данное соотношение


в виде линейной системы (18)


\begin{equation}


\begin{cases}


a_{11}g_1+a_{12}g_2+a_{13}g_3 = b_1,\\


a_{21}g_1+a_{22}g_2+a_{23}g_3 = b_2,\\


a_{31}g_1+a_{32}g_2+a_{33}g_3 = b_3,


\end{cases}


\end{equation}


где в соответствии с (16) и (17)


\begin{gather}


a_{ij}=\psi_{\delta(t_i)}^j(t_i),\ \ i,j=\overline{1,3};\\


b_i=x_i-\psi_{\delta(t_i)}^0(t_i)x_0,\ \ i=\overline{1,3}.


\end{gather}


Если система (33) совместна, то, разрешая ее относительно


коэффициентов $g_1$, $g_2$ и $g_3$, находим


решение обратной задачи (22), (23).





\begin{exam}


Рассмотрим сначала известное


дифференциально-разностное уравнение, когда


в (22) $a_0=1$, а $a_1=0$, т.\,е.


\begin{equation}


\Dot{x}(t)=x(t-1),\ \ t\in J=[0,3];\ \


x(t)=g(t),\ \ t\in J_0=[-1,0],


\end{equation}


с условиями


\begin{equation}


x(0)=x_0,\ \ x(t_1)=x_1,\ \ x(t_2)=x_2,\ \ x(t_3)=x_3.


\end{equation}


Согласно (32) и (27) определим функции $\psi_n^i(t)$ и


их значения в граничных точках. Результаты вычислений приведены


в табл.~1 и 2.


\medskip





\begin{center}


\hspace{100mm} Таблица 1 


\medskip





\begin{tabular}{|c|c|c|c|c|}


\hline \vrule height6mm width0mm depth3mm


%% \rule{0mm}{6mm}


$ i $ & $\psi_1^i(t)$ & $\psi_1^i(1)$ & $\psi_2^i(t)$ &


$\psi_2^i(2)$ \\


\hline \hline \vrule height6mm width0mm depth3mm


% \rule{0mm}{6mm}


$0$ & $1+t$ & 2 & 


$\frac{3}{2}+\frac{t^2}{2}$ & $\frac{7}{2}$ \\


\hline \vrule height6mm width0mm depth3mm


% \rule{0mm}{6mm} 


$1$ & $-t+\frac{t^2}{2}$ & $-\frac{1}{2}$ & 


$-\frac{7}{6}+\frac{3}{2}t-t^2+\frac{t^3}{6}$ & $-\frac{5}{6}$\\


%


\hline \vrule height6mm width0mm depth3mm


% \rule{0mm}{6mm}


$2$ & $t-t^2+\frac{t^3}{3}$ & $\frac{1}{3}$ & 


$\frac{5}{4}-\frac{7}{3}t+2t^2-\frac{2}{3}t^3+


\frac{t^4}{12}$ & $\frac{7}{12}$ \\


%


\hline \vrule height6mm width0mm depth3mm


% \rule{0mm}{6mm}


$3$ & $-t+\frac{3}{2}t^2-t^3+\frac{t^4}{4}$ & $-\frac{1}{4}$ & 


$-\frac{31}{20}+\frac{15}{4}t-4t^2+2t^3-\frac{t^4}{2}


+\frac{t^5}{20}$ & $-\frac{9}{20}$ \\


\hline


\end{tabular}


\end{center}


\medskip





\begin{center}


\hspace{73mm}Таблица 2


\medskip





\begin{tabular}{|c|c|c|}


\hline \vrule height6mm width0mm depth3mm


% \rule{0mm}{6mm}


$i$ & $\psi_3^i(t)$ & $\psi_3^i(3)$ \\


\hline


\hline \vrule height6mm width0mm depth3mm


% \rule{0mm}{6mm}


$0$ & $\frac{1}{6}+2t-\frac{t^2}{2}+\frac{t^3}{6}$


& $\frac{37}{6}$ \\


\hline \vrule height6mm width0mm depth3mm


% \rule{0mm}{6mm}


$1$ & $\frac{13}{6}-\frac{23}{6}t+2t^2-\frac{t^3}{2}+


\frac{t^4}{24}$ & $-\frac{35}{24}$\\


\hline \vrule height6mm width0mm depth3mm


% \rule{0mm}{6mm}


$2$ & $-\frac{197}{60}+\frac{19}{3}t-\frac{13}{3}t^2+


\frac{3}{2}t^3-\frac{t^4}{4}+\frac{t^5}{60}$ & $\frac{61}{60}$\\


\hline \vrule height6mm width0mm depth3mm


% \rule{0mm}{6mm}


$3$ & $\frac{331}{60}-\frac{237}{20}t+10t^2-\frac{9}{2}t^3+


\frac{9}{8}t^4-\frac{3}{20}t^5+\frac{t^6}{120}$ & $-\frac{47}{60}$\\


\hline


\end{tabular}


\end{center}





Пусть известны значения решения $x(t)$ в точках $t_i=i$,


$i=\overline{1,3}$. Тогда с учетом (34) и (35) перепишем (33)


в виде


\begin{equation}


\begin{cases}


\psi_1^1(1)g_1+\psi_1^2(1)g_2+\psi_1^3(1)g_3=


x_1-\psi_1^0(1)x_0,\\


\psi_2^1(2)g_1+\psi_2^2(2)g_2+\psi_2^3(2)g_3=


x_2-\psi_2^0(2)x_0,\\


\psi_3^1(3)g_1+\psi_3^2(3)g_2+\psi_3^3(3)g_3=x_3-\psi_3^0(3)x_0.


\end{cases}


\end{equation}


Подставляя сюда данные из табл.~1 и 2, приходим к линейной системе


\begin{equation} 


\begin{cases}


%%\left \{


%%\begin{array}{l}    %%% transform into \cases


-\frac{1}{2}g_1+\frac{1}{3}g_2-\frac{1}{4}g_3 = C_1,\\


-\frac{5}{6}g_1+\frac{7}{12}g_2-\frac{9}{20}g_3 = C_2,\\


-\frac{35}{24}g_1+\frac{61}{60}g_2-\frac{47}{60}g_3 = C_3,


%%\end{array}


%%\right.


\end{cases}


\end{equation}


где $C_1=x_1-2x_0$, $C_2=x_2-\frac{7}{2}x_0$,


$C_3=x_3-\frac{37}{6}x_0$.





Поскольку определитель этой системы $\det A=1$ отличен от нуля,


$\rank A=\rank \ov{A}=3$ и число линейно независимых опорных


функций $\psi_n^0(t)$ равно трем, задача (36), (37)


удовлетворяет условиям теоремы 1 и, следовательно, имеет


единственное решение. Разрешая (39) относительно неизвестных


коэффициентов $g_n$, получаем


\begin{align*}


g_1 &= 48x_1+600x_2-360x_3+24x_0, \\


g_2 &= 300x_1+2340x_2-1440x_3+90x_0, \\


g_3 &= 300x_1+1920x_2-1200x_3+80x_0.


\end{align*}


Так как согласно (23) и (24) $g_0=x_0$, данные соотношения


однозначно определяют начальную функцию


$g(t)=\sum\limits_{n=0}^3g_nt^n$,


порождающую решение $x(t)$ задачи (36), (37).


\medskip





{\bf Случай 1.} Вернемся к линейной системе (39) и рассмотрим


случай, когда в (37)


$$ 


x_0=1,\ \ x_1=2,\ \ x_2=\frac{7}{2},\ \ x_3=\frac{37}{6}.


$$


При этом $C_1=C_2=C_3=0$. Из (39) вытекает, что единственным


решением здесь будет тривиальное решение, т.\,е. $g_1=g_2=g_3=0$.


Тогда $g(t)=g_0=x_0=1$. Эта начальная функция порождает известное


решение начальной задачи (36), (37)


\begin{alignat*}{2}


x_1(t) &= 1+t, & \quad t&\in J_1=[0,1]; \\


x_2(t) &= 1+t+ \frac{(t-1)^2}{2}, &\quad t&\in J_2=[1,2]; \\


x_3(t) &= 1+t+ \frac{(t-1)^2}{2}+\frac{(t-2)^3}{6},


&\quad t&\in J_3=[2,3].


\end{alignat*}


Графики начальной функции $g(t)$ и этого решения приведены на рис.\,1


и обозначены сплошной линией.





\begin{center}


\unitlength=1mm


\begin{picture}(170,102) %% width, heigth, add 5 mm for signature


\put(51,100){\special {em:graph cherep.bmp}} %% left X, upper Y, -, /2, +toX


\put(78,0){\mbox Рис.\,1}


\end{picture}


\end{center}








Исследуем зависимость начальной функции от изменения одного из


данных, например, $x_1$.


\medskip





{\bf Случай 2.} Положим в (37)


$$ 


x_0=1,\ \ x_1=2{,}1,\ \ x_2=3{,}5,\ \ x_3=6{,}16666.


$$


Согласно (38) $g_1=4{,}8$, $g_2=g_3=30$ и


$$


g(t)=1+4{,}8t+30t^2+30t^3.


$$


Соответственно, для $x(t)$ на $J$ имеем


\begin{align*}


x_1(t)&= 1-3{,}8t+17{,}4t^2-20t^3+7{,}5t^4,\\


x_2(t)&= -13{,}1+49{,}7t-64{,}3t^2+40{,}8t^3-12{,}5t^4+1{,}5t^5,\\


x_3(t)&= 77{,}5666-181{,}9t+179{,}1t^2-92{,}2333t^3+26{,}45t^4-


4t^5+0{,}25t^6.


\end{align*}


Графики $g(t)$ и $x(t)$ обозначены на рис.\,1 коротким пунктиром.


\end{exam}





Исследуем обратную начальную задачу с условиями, наложенными не 


только на решение, но и на его производную.





\begin{exam}


Рассмотрим задачу (36) с условиями


равенства производных в точке $t=0$ у начальной функции и


порождаемого решения $x(t)$. Для этого перепишем (36) в виде


\begin{gather}


\Dot{x}(t)=x(t-1), \ \ t\in J=[0,1];\ \ 


x(t)=g(t),\ \ t\in J_0=[-1,0];\\


%


x(0)=x_0=1;\ \ x^{(n)}(0)=g^{(n)}(0)\ \ \text{при}\ \ t=0,\ \


n=\overline{1,3}.


\end{gather}


Так же, как в предыдущем случае, отметим, что здесь функции


$x(t)$ и $g(t)$ имеют производные до третьего порядка включительно.


При этом с учетом (24) $g(0)=


x_0=1$, $g^{(n)}(0)=n!g_n$, $n=\overline{1,3}$,


а производные $x_1^{(n)}(0)$ вычисляются согласно формулe (15) и


табл.\,1. Значения для данной задачи производных функций


$\psi_1^i(t)$, $i=\overline{0,3}$, при $t=0$, определяющих


$x_1^{(n)}(0)$, приведены в табл.\,3.


\medskip





\begin{center}


\hspace{43mm}Таблица 3


\medskip





\begin{tabular}{|c|c|c|c|}


\hline \vrule height6mm width0mm depth3mm


% \rule{0mm}{6mm}


$ i $ & $\frac{d}{dt}\psi_1^i(0)$ &


$\frac{d^2}{dt^2}\psi_1^i(0)$ &


$\frac{d^3}{dt^3}\psi_1^i(0)$ \\


\hline


\hline \vrule height6mm width0mm depth3mm


% \rule{0mm}{6mm}


$ 0 $ & 1 & 0 & 0 \\


\hline \vrule height6mm width0mm depth3mm


% \rule{0mm}{6mm}


$ 1 $ & $-1$  & 1 &  0 \\


\hline \vrule height6mm width0mm depth3mm


% \rule{0mm}{6mm}


$ 2 $ & 1 & $-2$  & 2 \\


\hline \vrule height6mm width0mm depth3mm


% \rule{0mm}{6mm}


$ 3 $ & $-1$  & 3 & $-6$  \\


\hline


\end{tabular}


\end{center}





Далее из условий (41) получаем


\begin{align*}


\frac{d}{dt}(\psi_1^1(0)g_1+\psi_1^2(0)g_2+\psi_1^3(0)g_3)&=


g_1-\frac{d}{dt}\psi_1^0(0), \\


\frac{d^2}{dt^2}(\psi_1^1(0)g_1+\psi_1^2(0)g_2+\psi_1^3(0)g_3)&=


2g_2- \frac{d^2}{dt^2}\psi_1^0(0), \\


\frac{d^3}{dt^3}(\psi_1^1(0)g_1+\psi_1^2(0)g_2+\psi_1^3(0)g_3)&=


6g_3- \frac{d^3}{dt^3}\psi_1^0(0).


\end{align*}





Подставляя данные из табл.\,3, приходим к следующей линейной


системе относительно неизвестных коэффициентов $g_n$:


$$


\begin{cases}


-2g_1+\enskip g_2-\enskip g_3=-1, \\


g_1-4g_2+3g_3= 0, \\


g_2-6g_3= 0.


\end{cases}


$$


Так как определитель этой системы $\det A=-45$, она однозначно


разрешима:


$$


g_1=0{,}567567,\ \ g_2=0{,}162162,\ \ g_3=0{,}027027.


$$


Тогда в силу теоремы 1 обратная начальная задача (40), (41)


имеет единственное решение в виде полинома третьей степени


\begin{equation}


g(t)=1+0{,}567567t+0{,}162162t^2+0{,}027027t^3,


\end{equation}


а порождаемое этой начальной функцией решение на основании


(25) и табл.\,1 представляется формулой


\begin{equation}


x_1(t)=\sum_{n=0}^{4}x_n^1t^n=


1+0{,}567567t+0{,}162162t^2+0{,}027027t^3+0{,}006757t^4.


\end{equation}


Отметим, что $x_n^1=g_n$, $n=\overline{0,3}$.





Графики функций $x_1(t)$ и $g(t)$ приведены на рис.\,1 


(длинный пунктир).


\end{exam}





Остановимся подробнее на результатах этого примера. Как известно,


начальная задача


$$


\Dot{x}(t)=x(t-1),\ \  x(0)=x_0=1,\ \


t\in J_{\infty}=(-\infty,\infty),


$$


имеет частное аналитическое решение, соответствующее вещественному


корню характеристического квазиполинома $k=e^{-k}$, который


получается при подстановке $x(t)=x_0e^{kt}$ в исходное уравнение.


В данном случае $k_1\approx 0{,}567143$. Представим это приближенное 


частное решение в виде ряда Маклорена


\begin{multline*}


\ov{x}(t)=e^{0{,}567143t}=\sum_{n=0}^{\infty}


\frac{(0{,}567143)^n}{n!}t^n= \\


%


=1+0{,}567143t+0{,}16082t^2+0{,}03040t^3+0{,}00431t^4+


\sum_{n=5}^{\infty}\frac{(0{,}567143)^n}{n!}t^n.


\end{multline*}


Сравнение коэффициентов данного ряда с коэффициентами в формулах


(42) и (43) позволяет сделать 


\medskip





{\bf Вывод.} Функции $g(t)$ и $x_1(t)$, определенные 


соответственно формулами (42), (43) и продолженные 


на $J_T=[-a,b]$, можно рассматривать как приближения к частному 


аналитическому решению задачи (41) на $J_T$.








\begin{thebibliography}{9}





\bibitem{}


Мышкис А.Д. {\em Линейные дифференциальные уравнения с 


запаздывающим аргументом.} -- М.--Л.: Гостехиздат, 1951. -- 256~с.


\bibitem{2}


Пинни Э. {\em Обыкновенные дифференциально-разностные уравнения.} -- 


М.: Ин. лит., 1961. -- 248~с.


\bibitem{3}


Беллман Р., Кук К.Л. {\em Дифференциально-разностные уравнения.} 


-- М.: Мир, 1967. -- 548~с.


\bibitem{4}


Хейл Д. {\em Теория функционально-дифференциальных уравнений.} -- 


М.: Мир, 1984. -- 411~с.


\bibitem{5}


Азбелев Н.В., Максимов В.П, Рахматуллина Л.Ф. {\em Введение в теорию


функционально-дифференциальных уравнений.} -- М.: Наука, 1991. -- 280~с.


\bibitem{6}


Baker C.T.H., Bocharov G.A., Paul C.A.H., Rihan F.A. {\em Modelling and


analysis of time-lags in some basic patterns of cell proliferation}


// J. Math. Biol. -- 1998. -- V.\,37. -- P.\,341--371.


\bibitem{7}


Габасов З., Кириллова Ф.М. {\em Качественная теория оптимальных 


процессов.} -- М.: Наука, 1971. -- 507~с. 


\bibitem{8}


Забелло Л.Е., Копейкина Т.Б. {\em Управляемость по начальной функции


систем с запаздыванием} // Дифференц. уравнения. -- 1976. -- 


Т.\,12. -- \No\,12. -- С.\,2267--2268. 





\end{thebibliography}





\vskip 2cm





\post{\it Институт динамики систем}{\it Поступила}


\post{\it и теории управления Сибирского}{28.11.2002}


\post{\it отделения Российской академии наук}{}


\post{\it }{}





\label{end}


\end{document}


